% =============================================================================
% БИЛЕТ 2
% =============================================================================

\ticket{2}{}

\question{Дифференцируемость ФНП. Частные производные, производные по направлению, матрица Якоби}

\begin{defn}[Дифференцируемость ФНП в точке]
  Функция $f : \mathbb{R}^n \to \mathbb{R}$ называется \textit{дифференцируемой}
  в точке $x^0$, если существует линейная функция $L : \mathbb{R}^n \to \mathbb{R}$
  такая, что:
  \[
    f(x) - f(x^0) = L(x - x^0) + o(\|x - x^0\|), \quad x \to x^0.
  \]
  Линейную функцию $L$ называют \textit{дифференциалом} $f$ в точке $x^0$
  и обозначают $df(x^0)$. Так как любое линейное отображение
  $\mathbb{R}^n \to \mathbb{R}$ задаётся скалярным произведением, то:
  \[
    f(x) - f(x^0) = \sum_{i=1}^n A_i(x_i - x_i^0) + o(\|x - x^0\|),
  \]
  где вектор $\nabla f(x^0) := (A_1, \ldots, A_n)$ называется
  \textit{градиентом} $f$ в точке $x^0$.
\end{defn}

\begin{defn}[Частная производная]
  Пусть $f : \mathbb{R}^n \to \mathbb{R}$, $x^0 \in \mathbb{R}^n$.
  \textit{Частной производной} $f$ по переменной $x_k$ в точке $x^0$
  называется предел:
  \[
    \frac{\partial f}{\partial x_k}(x^0)
    := \lim_{t \to 0}
       \frac{f(x^0 + t e_k) - f(x^0)}{t},
  \]
  если он существует, где $e_k = (0,\ldots,\underset{k}{1},\ldots,0)$~---
  орт $k$-й координатной оси.
\end{defn}

\begin{defn}[Производная по вектору]
  Пусть $f : \mathbb{R}^n \to \mathbb{R}$, $x^0 \in \mathbb{R}^n$,
  $l \in \mathbb{R}^n \setminus \{0\}$.
  \textit{Производной $f$ по вектору $l$} в точке $x^0$ называется предел:
  \[
    \frac{\partial f}{\partial l}(x^0)
    := \lim_{t \to +0} \frac{f(x^0 + tl) - f(x^0)}{t},
  \]
  если он существует. Если $\|l\| = 1$, то говорят о
  \textit{производной по направлению}.
\end{defn}

\begin{theorem}[Первое НУ дифференцируемости]
  Если $f : \mathbb{R}^n \to \mathbb{R}$ дифференцируема в точке $x^0$,
  то она непрерывна в этой точке.
  \begin{proof}
    Из определения дифференцируемости:
    \[
      f(x) - f(x^0)
      = \sum_{i=1}^n A_i(x_i - x_i^0) + o(\|x - x^0\|).
    \]
    При $x \to x^0$ каждое слагаемое $A_i(x_i - x_i^0) \to 0$
    и $o(\|x - x^0\|) \to 0$, значит $f(x) \to f(x^0)$.
  \end{proof}
\end{theorem}

\begin{theorem}[Второе НУ: формула для производной по направлению]
  Если $f : \mathbb{R}^n \to \mathbb{R}$ дифференцируема в точке $x^0$,
  то для любого $l \in \mathbb{R}^n \setminus \{0\}$ существует производная
  по вектору $l$ и:
  \[
    \frac{\partial f}{\partial l}(x^0)
    = df(x^0)(l)
    = \langle \nabla f(x^0),\, l \rangle
    = \sum_{i=1}^n A_i l_i.
  \]
  \begin{proof}
    Подставим $x = x^0 + tl$ в определение дифференцируемости:
    \[
      f(x^0 + tl) - f(x^0)
      = t \sum_{i=1}^n A_i l_i + \varepsilon(t)\cdot t\|l\|,
      \quad \varepsilon(t) \to 0.
    \]
    Делим на $t > 0$ и переходим к пределу при $t \to +0$:
    \[
      \frac{f(x^0 + tl) - f(x^0)}{t}
      = \sum_{i=1}^n A_i l_i + \varepsilon(t)\cdot\|l\|
      \;\xrightarrow{t \to +0}\;
      \sum_{i=1}^n A_i l_i. \qedhere
    \]
  \end{proof}
\end{theorem}

\begin{corollary}[Третье НУ: существование частных производных]
  Если $f$ дифференцируема в $x^0$, то для каждого $k \in \{1,\ldots,n\}$
  существует частная производная:
  \[
    \frac{\partial f}{\partial x_k}(x^0) = A_k = \bigl(\nabla f(x^0)\bigr)_k.
  \]
  \begin{proof}
    Подставим $l = e_k$ во второе НУ:
    \[
      \frac{\partial f}{\partial x_k}(x^0)
      = \frac{\partial f}{\partial e_k}(x^0)
      = \sum_{i=1}^n A_i (e_k)_i = A_k. \qedhere
    \]
  \end{proof}
\end{corollary}

\begin{defn}[Матрица Якоби]
  Пусть $F = (F_1, \ldots, F_m) : \mathbb{R}^n \to \mathbb{R}^m$
  дифференцируемо в точке $x^0$.
  \textit{Матрицей Якоби} отображения $F$ в точке $x^0$ называется
  матрица размера $m \times n$:
  \[
    \mathcal{D}F(x^0)
    :=
    \begin{pmatrix}
      \dfrac{\partial F_1}{\partial x_1}(x^0) & \cdots &
      \dfrac{\partial F_1}{\partial x_n}(x^0) \\[8pt]
      \vdots & \ddots & \vdots \\[4pt]
      \dfrac{\partial F_m}{\partial x_1}(x^0) & \cdots &
      \dfrac{\partial F_m}{\partial x_n}(x^0)
    \end{pmatrix}.
  \]
  Для функции $f : \mathbb{R}^n \to \mathbb{R}$ матрица Якоби~---
  это строка-градиент $\nabla f(x^0)^\top$.
\end{defn}
\question{Непрерывность суммы равномерно сходящегося ряда из непрерывных функций. Почленное интегрирование/дифференцирование ФР}

% Ответ...

